\section{K-014: the complete sixteenth-root pole in dimension five}

\begin{theorem}[K-014]
Let $\zeta=e^{2\pi i/16}$.  The fixed scheme of the grading group
$\mu_{16}$ on
$X_5=\operatorname{Spec}\mathbb C[(\mathbb C^5)^{\otimes3}]^{SL_5^3}$
has exactly $424$ irreducible curve components.  Every nonvertex
point lies on their nonzero grading orbits, where the quotient and
the fixed curve are smooth.  There are no other positive-dimensional
components.  Moreover
\[
 \boxed{\displaystyle
 \lim_{t\to\zeta}(1-t^{16})H^{\rm ten}_5(t)
 =\frac{78+446\zeta-211\zeta^2-301\zeta^3-301\zeta^4
              -211\zeta^5+446\zeta^6+78\zeta^7}{32768}\ne0.}
                                                               \tag{K-014a}
\]
Thus the Hilbert series has a simple pole at every primitive
sixteenth root of unity.
\end{theorem}

\begin{proof}
K-011 covers every nonzero grading-fixed quotient point by one of
the seventy-seven spectral types, allowing every possible
two-primary lift up to the bound supplied by K-009.  K-010 shows
that each fixed space contains a single product-linear orbit of
semistable tensors.  K-012 and K-013 show that the full stabilizer
of every representative is $S\times\mu_{16}$, with faithful
determinant character.  Therefore the effective $SL_5^3$ stabilizer
is trivial and the quotient is smooth there by the smooth slice
description.

In the seventy-six simple-spectrum cases the fixed tensor space
has dimension $13$ and the fixed determinant-one Lie algebra has
dimension $15-3=12$.  In the repeated-spectrum case these dimensions
are $15$ and $17-3=14$.  The fixed quotient tangent dimension is one
in both cases.  The grading orbit has stabilizer $\mu_{16}$, so its
one-dimensional closure is the entire local fixed curve component.
Every nonzero fixed point is on one of these orbits.  As the fixed
scheme of a finite-order automorphism on a smooth complex variety
is smooth, it is generically reduced along all these curves.

We must restore the order of the three tensor factors.  For each
spectral type consider all six factor permutations and identify
those related by coordinate permutations and scalar shifts whose
sum is zero.  Supported-edge shifts provide the complete
normalization, as in K-010.  Sixty-five types yield six ordered
classes, eleven yield three, and one yields one.  Consequently
\[
                        65\cdot6+11\cdot3+1=424.             \tag{K-014b}
\]
These ordered classes give distinct product-linear orbits.
Indeed, the full stabilizer modulo scalars has exactly one element
with determinant character $\zeta$.  A product-linear equivalence
of two tensors would identify those elements, forcing their ordered
spectra to agree up to coordinate permutations and the already
allowed scalar shifts.  Thus the count does not overcount fixed
curves.

Every curve has leading density one.  Its restriction algebra on
the grading orbit is a finitely generated semigroup algebra
$\mathbb C[u^s:s\in\Gamma]$.  The exact orbit stabilizer implies
$\gcd\Gamma=16$.  The numerical semigroup $\Gamma/16$ is cofinite,
so the curve has Hilbert function eventually one at multiples of
$16$ and zero elsewhere.  Thus its normalized leading density is
$e_Y=1$, exactly as in the rechecked P-F60o; possible initial holes
at the vertex do not alter this density.

For a spectral triple $W=(A,B,C)$, the moving normal determinant
coefficient is
\[
 C_W(\zeta)=
 \frac{\displaystyle
     \prod_{X\in\{A,B,C\}}
       \prod_{a,b\in X:\ a-b\not\equiv0\ (16)}
                           (1-\zeta^{a-b})}
      {\displaystyle
       \prod_{(a,b,c)\in A\times B\times C:\ a+b+c\not\equiv0\ (16)}
                           (1-\zeta^{a+b+c})}.              \tag{K-014c}
\]
Multisets and their multiplicities are used in both products.  The
finite stabilizer has now been proved trivial in the effective
determinant-one action, so this is the actual local quotient
coefficient, not merely a determinant before an unknown finite
quotient.  Its convention is the point pullback
$F\mapsto F\circ g$, whose eigenvalue in degree $s$ is $\zeta^s$;
on cotangents it has the same eigenvalues as the point differential,
as explicitly checked in P-F60p.  Subtracting the moving Lie
representation from the moving tensor representation gives
(K-014c).

The exact script \texttt{n5\_residue.py} multiplies these factors in
$\mathbb Q[z]/(z^8+1)$, using only integers and rational fractions,
and inverts the nonzero denominators by an eight-dimensional rational
linear solve.  It also recomputes all ordered-factor multiplicities.
The result, as coefficients in the basis $1,z,\ldots,z^7$, is
\[
       \sum_W m_W C_W(z)
       =\frac1{32768}(78,446,-211,-301,-301,-211,446,78).
                                                               \tag{K-014d}
\]
All seventy-seven individual products and the multiplicities are
recorded in \nolinkurl{n5_residue_verified.json}.

The multiple-component formula P-F23g now adds these $424$ generic
curve contributions.  Lower-dimensional strata, all supported at
the vertex, contribute no simple-pole coefficient.  Equations
(K-014c)--(K-014d) prove (K-014a).  Since
$1,\zeta,\ldots,\zeta^7$ are rationally independent and the displayed
coefficient vector is nonzero, this residue is nonzero.
Galois conjugation proves nonvanishing at each primitive sixteenth
root, and the one-dimensional fixed locus gives simple pole order.
\end{proof}

The strictly proper part of the Hilbert series supported at these
primitive-root poles is exactly
\[
 \frac{78+446t-211t^2-301t^3-301t^4-211t^5+446t^6+78t^7}
      {65536(1+t^8)}.                                       \tag{K-014e}
\]
Indeed $(1-t^{16})/(1+t^8)=1-t^8$ takes the value $2$ at every
primitive sixteenth root.  Its coefficient sequence is the
constant-constituent periodic wave
\[
 \frac1{65536}
 (78,446,-211,-301,-301,-211,446,78,
  -78,-446,211,301,301,211,-446,-78),                        \tag{K-014f}
\]
of period sixteen.  These exact finite-dimensional results extend
the calibration through $n=5$.  They do not by themselves decide
the existence of primitive $2^{n-1}$ poles for every $n$, nor
polynomial bounds on the positivity and saturation indices.
